% !TEX root = ../main.tex
\section{\PCER}
\label{sec:method}

\subsection{기호}

$T_0$을 순위의 동결 시점, $C$를 $T_0$까지 Stacks 메인넷에 배포된 컨트랙트의 집합, $\mathcal{W}=(T_0-\Delta,T_0]$를 관측 창이라 하자. ERC-20 연구와 같이 $\Delta=180$일로 둔다 \cite{kwon2026endorserank}. $\mathcal{W}$ 안의 트랜잭션 $\tau$에 대해 $o(\tau)$는 발신자(standard principal이며, 아래에서는 지갑이라 부른다), $c(\tau)\in C$는 페이로드가 컨트랙트 호출일 때 페이로드에 적힌 컨트랙트, $f(\tau)$는 수수료, $\mathrm{out}_a(\tau)$는 $\tau$ 동안 $o(\tau)$가 보내거나 소각한 자산 $a$의 양으로, 트랜잭션의 자산 이벤트에서 읽는다. $\mathrm{Auth}(\tau)$는 \ref{sec:mapping-rules}절의 규칙에 따라 $\tau$가 승인 포스트컨디션을 담고 있는 자산들의 집합이며, Allow 모드의 포괄 승인도 포함한다. $p(a)$는 공통 단위로 나타낸 자산 $a$의 $T_0$ 시점 가격이다. 믿을 만한 가격이 없는 자산과, 주 분석에서는 대체 불가능 토큰에 대해 $p(a)=0$이다.

\subsection{승인 기록과 최신 상태 규칙}
\label{sec:method-records}

\begin{definition}[승인 기록]
지갑 $u$, 컨트랙트 $c$, 자산 $a$에 대해 $\tau^{*}(u,c,a)$를 $\mathcal{W}$ 안에서 $o(\tau)=u$, $c(\tau)=c$, $a\in\mathrm{Auth}(\tau)$를 만족하는 가장 늦은 성공 트랜잭션이라 하자. 그 값은 $v(u,c,a)=p(a)\,\mathrm{out}_a(\tau^{*})$이다. $u$의 $c$에 대한 승인은 $A(u,c)=\sum_a v(u,c,a)$, $u$가 승인한 컨트랙트의 집합은 $C_u=\{c : \text{어떤 } a\text{에 대해 } \tau^{*}(u,c,a)\text{가 존재}\}$이고, $u$의 보증에서 $c$가 차지하는 몫은 다음과 같다.
\begin{equation}
\label{eq:share}
s_u(c)=
\begin{cases}
A(u,c)\big/\sum_{c'\in C_u}A(u,c') & \text{분모가 양수일 때,}\\
1/|C_u| & \text{그 밖의 경우.}
\end{cases}
\end{equation}
\end{definition}

두 가지 선택을 설명해 둔다. 첫째, allowance 값 자체를 쓰는 EndorseRank와 달리, 값은 포스트컨디션에 적힌 한도가 아니라 지갑에서 실제로 나간 양이다. 성공한 트랜잭션에서 상한은 유출보다 작을 수 없고, 한도는 ERC-20 allowance처럼 지갑이 가진 양보다 훨씬 크게 비용 없이 적을 수 있다. 실제 유출량은 자산이 움직여야 생긴다. 둘째, 최신 상태 규칙은 EndorseRank의 규칙 그대로이며, 소유자·지출자·토큰 대신 지갑·컨트랙트·자산마다 적용한다. 같은 컨트랙트를 천 번 호출한 지갑도 자기 몫 배분에는 가장 최근 호출의 값으로 한 번만 들어간다. 반복 호출은 그 호출에 든 수수료를 통해서만 지갑의 가중치를 높인다(\ref{sec:method-weights}절).

\subsection{재시작을 바꿔야 하는 이유}

EndorseRank를 더 손대지 않고 Stacks로 옮기면, 지갑과 컨트랙트가 노드이고 간선 $u\to c$를 $A(u,c)$로 가중하며 컨트랙트 사이에도 간선이 있는 그래프에 균등 재시작을 쓰게 된다. 다음 명제는 이런 그래프가 무엇을 재는지 보여 준다.

\begin{proposition}[균등 재시작은 지갑 수를 센다]
\label{prop:bipartite}
$G$의 노드 집합이 $N$개 노드로 된 $W\cup C$이고, 지갑에서 컨트랙트로 가는 간선 $u\to c$의 전이 확률이 $s_u(c)$이며, 컨트랙트 사이에도 간선이 있고, 재시작이 균등한 $\pi_v=1/N$이라 하자. 그러면 모든 지갑의 점수는 $r_W=(1-d+d\,\delta)/N$으로 같다. 여기서 $\delta$는 댕글링 노드가 가진 점수이다. 또 지갑들이 컨트랙트 $c$에 넘기는 점수는 $d\,r_W\sum_{u}s_u(c)$이다. 지갑 가운데 $k$개가 저마다 컨트랙트 $x$만 승인하면, 이 지갑들은 승인한 양과 관계없이 모두 합쳐 $k\,d\,r_W$를 $x$에 넘긴다. 이때 $N$, $\delta$, 따라서 $r_W$는 그 지갑들을 포함한 그래프의 값이다.
\end{proposition}

\begin{proof}
지갑으로 들어오는 간선은 없다. 컨트랙트는 트랜잭션을 발신하지 않고, 컨트랙트 사이의 간선은 컨트랙트에서 끝나기 때문이다. 따라서 식~\eqref{eq:pagerank}에 의해 각 지갑의 점수는 $(1-d)/N+d\,\delta/N$이며 지갑에 따라 달라지지 않는다. 지갑 $u$는 $d\,r_W$를 나가는 간선을 따라 $s_u(c)$의 비율로 넘기므로, $c$에 넘어가는 양의 합은 $d\,r_W\sum_u s_u(c)$이다. $x$만 승인한 지갑은 $s_u(x)=1$이다.
\end{proof}

균등 재시작에서 승인한 양은 각 지갑이 정해진 기여분을 자기 컨트랙트들에 어떻게 나누는지만 정한다. 컨트랙트 점수 가운데 지갑에서 오는 부분은, 지갑마다 한 표를 자기가 쓰는 컨트랙트들에 나누어 준 뒤 센 지갑 수이다. 시빌 지갑 하나는 트랜잭션 수수료 한 번으로 온전한 한 표를 가져온다. 금액으로 가중해도 이것은 바뀌지 않는다. 이 명제는 ERC-20 연구에서 보류해 둔 승인과 EndorseRank보다 단순 입차수가 더 잘 맞은 까닭도 하나 시사한다 \cite{kwon2026endorserank}. 지출자이면서 소유자이기도 한 주소가 적은 allowance 그래프는 이분 그래프에 가깝고, 그런 그래프에서 걷기는 개수에 보태는 것이 거의 없다.

\subsection{지갑 가중치와 재시작 분포}
\label{sec:method-weights}

그래서 지갑을 노드 집합에서 빼고, 공격자가 돈을 내야 얻을 수 있는 가중치 $\omega(u)$를 붙여 재시작 분포로만 작용하게 한다.
\begin{equation}
\label{eq:restart}
\omega(u)=\sum_{\tau\in\mathcal{W}:\,o(\tau)=u}\min\bigl(f(\tau),f_{\max}\bigr),
\qquad
\pi_c=\frac{\sum_{u}\omega(u)\,s_u(c)}{\sum_{u:\,C_u\neq\emptyset}\omega(u)}
\end{equation}
지갑의 가중치는 그 지갑이 창 안에서 발신한 모든 트랜잭션의 수수료를 합한 값이며, 각 수수료는 $\mathcal{W}$ 안 수수료의 99번째 백분위수인 $f_{\max}$에서 자른다. 대납자가 낸 수수료도 발신자의 것으로 센다. 누군가는 그 수수료를 냈고, 자기 지갑의 수수료를 대납하는 공격자는 그 돈을 스스로 낸다. 식~\eqref{eq:restart}는 $\omega$에 비례해 고른 지갑에서 출발해 승인 간선을 한 걸음 따라간 걷기가 보게 되는 재시작 분포와 같다.

\begin{proposition}[수수료 예산을 나눠도 얻는 것이 없다]
\label{prop:split}
공격자가 $x$만 승인하는 지갑들을 통해 수수료 예산 $F$를 쓴다고 하자. 식~\eqref{eq:restart}에서 이 지갑들이 $\pi_x$의 분자에 보태는 양은 예산을 지갑들에 어떻게 나누든 $F$ 이하이다. 반면 $g(0)=0$인 오목 함수 $g$로 $\omega(u)=g(u\text{의 수수료})$를 쓰면, $F$를 $k$개 지갑에 고르게 나눌 때 $k\,g(F/k)\geq g(F)$를 보탠다. $\omega\equiv 1$이면 최소 수수료 $k$번의 비용으로 $k$를 보탠다.
\end{proposition}

\begin{proof}
첫 주장은 선형성과, 센 수수료를 낮출 수만 있는 상한에서 나온다. $g(0)=0$인 오목 함수 $g$는 준가법적이므로 $g(F)=g\bigl(\sum_{i=1}^{k}F/k\bigr)\leq\sum_{i=1}^{k}g(F/k)$이다.
\end{proof}

선형 가중치에서는 영향력이 트랜잭션에 쓴 돈에 비례하고, 공격자는 신원을 늘려서 이 비용을 피할 수 없다. 그 대가로 거래를 많이 하는 지갑이 드물게 하는 지갑보다 크게 셈해진다. 지갑 하나에 한 표를 주는 규칙에는 Stacks에 없는 신원 계층이 필요하며, Cheng과 Friedman의 결과에 따르면 대칭적인 규칙으로는 이 문제를 피할 수 없다 \cite{cheng2005sybilproof}. $\omega\equiv1$에 맞서는 공격자는 최소 수수료 한 번으로 영향력 한 단위를 산다. 식~\eqref{eq:restart}에서는 정직한 지갑 $n$개의 영향력을 따라잡으려면 그 지갑들이 낸 만큼을 내야 하므로, 비용은 정직한 지갑의 평균 수수료와 최소 수수료의 비만큼 오른다. \ref{sec:evaluation}절에서 이 비를 메인넷에서 잰다. 민감도 분석에서는 $\mathcal{W}$ 동안 지갑의 STX 잔액을 시간 평균한 자본-일도 쓴다. 자본-일은 수수료처럼 나눠도 합이 변하지 않고, 같은 코인이 동시에 두 지갑에 있을 수도 없다.

Stacks 채굴자는 트랜잭션 수수료를 받으므로, 채굴자는 수수료를 자기 자신에게 낼 수 있다. 상한 $f_{\max}$는 스스로 낸 트랜잭션 하나가 살 수 있는 양을 제한한다. 이 문제는 \ref{sec:discussion}절에서 다시 다룬다.

\subsection{컨트랙트 사이의 간선}

$E_{\mathrm{del}}$과 $E_{\mathrm{dep}}$을 \ref{sec:mapping}절의 위임 간선과 의존 간선이라 하자. 이 간선들은 소스와 거버넌스 트랜잭션에서 다음과 같이 뽑는다.

\begin{itemize}
\item \textbf{의존:} $A$에 대상 $B$를 절대 principal, 상대 principal \code{.name}(같은 배포자의 컨트랙트), 그런 principal에 묶인 상수(Clarity~5) 중 하나로 쓴 \code{contract-call?}이 있거나, $A$의 \code{impl-trait} 또는 \code{use-trait}이 $B$에 정의된 trait을 가리키면 $A\to B$이다.
\item \textbf{위임:} 대상을 알아낸 $A$에서 $B$로의 호출이 \code{as-contract?}, 옛 \code{as-contract}, \code{restrict-assets?}의 본문 안에 있거나, $A$가 ExecutorDAO 패턴을 따르고 $A$의 이력에서 가장 최근의 \code{set-extension} 호출이 $B$를 켜 두었으면 $A\to B$이다.
\item \textbf{제외:} 대상이 trait 타입 인자인 호출, 자기 자신으로 가는 간선, $T_0$에 존재하지 않는 컨트랙트에 대한 참조는 뺀다.
\end{itemize}

두 집합에 모두 든 간선은 위임으로 센다. 가중치는 $\beta_{\mathrm{del}}=2$, $\beta_{\mathrm{dep}}=1$로 미리 정하고, 민감도 분석에서 $\beta_{\mathrm{del}}$을 $\{1,2,4\}$로 바꾼다. 위임을 더 무겁게 두는 것은 위임이 자산이나 권한을 피호출자에게 맡기기 때문이며, 승인과 사용을 구분하는 것이 EndorseRank의 바탕이다. 전이 행렬은 $P_{AB}=\beta(A,B)/\sum_{B'}\beta(A,B')$이다.

\subsection{점수}

\PCER는 노드 집합 $C$, 전이 행렬 $P$, 식~\eqref{eq:restart}의 재시작 $\boldsymbol{\pi}$로 식~\eqref{eq:pagerank}를 푼 해이다. ERC-20 연구와 같이 $d=0.85$, $\ell_1$ 허용 오차 $10^{-8}$, 최대 300회 반복을 쓴다 \cite{kwon2026endorserank}. 알고리즘~\ref{alg:pcer}에 단계를 적었다. 점수가 0인 컨트랙트는 사용된 근거가 없는 컨트랙트이며, 레지스트리는 이들을 목록에는 올리되 순위는 매기지 않는다.

\begin{proposition}[공격자 컨트랙트의 점수]
\label{prop:region}
집합 $S\subseteq C$에 대해 $\Phi_{\mathrm{in}}(S)=\sum_{v\notin S}r_v P_{vS}$를 $S$ 밖의 컨트랙트에서 간선을 따라 $S$로 흘러드는 점수라 하자. 여기서 $P_{vS}=\sum_{y\in S}P_{vy}$이다. 그러면 다음이 성립한다.
\begin{equation}
r(S)\;\leq\;\frac{\pi(S)}{1-d}+\frac{d}{1-d}\,\Phi_{\mathrm{in}}(S)
\end{equation}
특히 $\pi(S)=0$이고 $S$ 밖의 어떤 컨트랙트도 $S$로 가는 간선을 갖지 않으면 $r(S)=0$이다.
\end{proposition}

\begin{proof}
식~\eqref{eq:pagerank}를 $S$에 대해 더하면 $r(S)=(1-d)\pi(S)+d\sum_{v}r_vP_{vS}+d\,\delta\,\pi(S)$이다. 나가는 간선이 있는 $S$의 노드는 많아야 자기 점수 전부를 $S$로 보내고, 댕글링 노드는 아무것도 보내지 않으므로, $S$ 안의 댕글링 점수를 $\delta_S$라 하면 $\sum_{v\in S}r_vP_{vS}\leq r(S)-\delta_S$이다. 따라서 $(1-d)\,r(S)\leq(1-d)\pi(S)+d\,\Phi_{\mathrm{in}}(S)+d\,(\delta\,\pi(S)-\delta_S)$이다. $\delta\leq1$이고 $\delta_S\geq0$이므로 마지막 항은 $d\,\pi(S)$ 이하이며, 여기서 상한이 나온다. $\pi(S)=0$이고 $\Phi_{\mathrm{in}}(S)=0$이면 부등식은 $r(S)\leq 0$이 된다.
\end{proof}

우변의 두 항은 공격자가 얻어야 하는 두 가지이다. $\pi(S)$는 $S$ 안의 컨트랙트를 승인하는 지갑들이 낸 수수료로 사며, 그 값은 명제~\ref{prop:split}가 정한다. $\Phi_{\mathrm{in}}(S)$를 얻으려면 이미 점수가 있는 컨트랙트가 코드에서 $S$를 참조하거나 $S$를 확장으로 켜야 한다. 이것은 다른 개발자와 거버넌스가 내리는 결정이며, 소셜 그래프 방어에서 말하는 공격 간선에 해당한다 \cite{yu2006sybilguard,cao2012sybilrank}. 대상을 참조하는 컨트랙트를 배포하면 $S$의 구성원은 늘지만 어느 항도 커지지 않는다. 반면 균등 재시작에서는 배포한 컨트랙트마다 재시작 질량 $1/|C|$를 함께 가져오므로, 공격자가 배포한 컨트랙트 수만큼 상한이 커진다.

\begin{algorithm}[htbp]
\caption{\PCER}
\label{alg:pcer}
\begin{algorithmic}[1]
\Require $T_0$ 시점의 소스가 있는 컨트랙트 $C$; $\mathcal{W}$ 안의 트랜잭션; 가격 $p$; 매개변수 $d,\beta_{\mathrm{del}},\beta_{\mathrm{dep}},f_{\max}$
\Ensure $C$ 위의 점수 $\mathbf{r}$
\For{$\mathcal{W}$ 안의 성공한 컨트랙트 호출 $\tau$와 각 $a\in\mathrm{Auth}(\tau)$}
  \State $\tau$가 저장된 것보다 늦으면 $\tau^{*}(o(\tau),c(\tau),a)$로 저장
\EndFor
\State $A(u,c)\gets\sum_a p(a)\,\mathrm{out}_a(\tau^{*}(u,c,a))$; 식~\eqref{eq:share}로 $s_u(c)$ 계산
\State $\omega(u)\gets\sum_{\tau:\,o(\tau)=u}\min(f(\tau),f_{\max})$; 식~\eqref{eq:restart}로 $\boldsymbol{\pi}$ 계산
\State 각 컨트랙트를 파싱; 리터럴·상대·상수 대상을 풀어 $E_{\mathrm{dep}}$과 $E_{\mathrm{del}}$을 모음; 켜진 확장을 $E_{\mathrm{del}}$에 더함
\State $P_{AB}\gets\beta(A,B)/\sum_{B'}\beta(A,B')$
\State $\mathbf{r}\gets\boldsymbol{\pi}$
\Repeat
  \State $\mathbf{r}'\gets(1-d)\boldsymbol{\pi}+dP^{\top}\mathbf{r}+d\,\delta(\mathbf{r})\,\boldsymbol{\pi}$; $\mathbf{r}$과 $\mathbf{r}'$를 맞바꿈
\Until{$\lVert\mathbf{r}-\mathbf{r}'\rVert_1<10^{-8}$ 또는 300회 반복}
\State \Return $\mathbf{r}$
\end{algorithmic}
\end{algorithm}

\subsection{합성 예제}
\label{sec:method-toy}

명제들이 실제로 어떻게 작동하는지 보이려고 작은 합성 그래프를 만들었다(스크립트는 논문 저장소에 있고, 매개변수는 부록~\ref{app:toy}에 적었다). 정직한 컨트랙트 \ToyContracts 개에는 SIP-010 trait 하나, 토큰 여섯 개, \code{as-contract?} 안에서 토큰을 내주는 풀 세 개, 라우터, 볼트, DAO가 들어 있다. 정직한 지갑 \ToyWallets 개는 저마다 진입 컨트랙트를 한 개에서 세 개 승인하고, 창 동안 평균 최소 수수료의 \ToyMeanFee 배를 낸다. 공격자는 토큰 $x$를 배포하고 \ref{sec:evaluation}절의 공격 가운데 세 가지로 $x$를 상위 \ToyTop 위 안에 올리려 한다. A1은 $x$를 호출하는 컨트랙트를 배포한다(컨트랙트 하나에 수수료 \ToyDeployFee 단위). A2는 최소 수수료로 $x$를 한 번 승인하는 지갑을 만든다. A4는 $x$를 trait 인자로 넣어 정직한 라우터를 호출한다. 왕복 송금인 A3은 AWP에만 통하므로 이 예제에서는 뺐다. 표~\ref{tab:toy-cost}에 다섯 가지 순위 규칙에서 상위 \ToyTop 위에 드는 데 필요한 최소 비용을 적었다.

\begin{table}[htbp]
\centering
\caption{정직한 컨트랙트 \ToyContracts 개로 된 합성 그래프에서 공격자 토큰을 상위 \ToyTop 위 안에 올리는 최소 공격 비용(최소 호출 수수료 단위). 공격 크기는 두 배씩 늘려 시험한 뒤 이분 탐색으로 좁히며, 공격이 커져도 토큰의 순위가 내려가지 않는다고 가정한다. 동점은 토큰에 유리하게 센다. ``도달 못 함''은 공격 크기가 200,000일 때도 토큰이 상위 \ToyTop 위 밖에 있다는 뜻이다. 그래프는 꾸며 낸 것이며, 이 표는 명제~\ref{prop:bipartite}--\ref{prop:region}를 예시할 뿐 Stacks 메인넷에 대한 결과가 아니다.}
\label{tab:toy-cost}
\small
% Generated by scripts/toy_sybil_example.py. Do not edit by hand.
\begin{tabular}{@{}lrrr@{}}
\toprule
Scoring rule & A1 & A2 & A3 \\
\midrule
Distinct callers & not reached & 59 & not reached \\
Dependency PageRank, uniform restart & 10 & not reached & not reached \\
EndorseRank on Stacks, uniform restart & 750 & 75 & not reached \\
PC-EndorseRank (fee-weighted restart) & not reached & 2,234 & not reached \\
PC-EndorseRank with runtime trait bindings & not reached & 2,234 & 3,952 \\
\bottomrule
\end{tabular}

\end{table}

기준선은 저마다 적어도 한 가지 공격에 무너진다. 호출자 수 세기는 지갑 농장(A2)에 \ToyCostCallersATwo 단위로 무너진다. teaRank와 같은 꼴인 균등 재시작 의존성 PageRank는 농장 컨트랙트 하나(A1, \ToyCostDepPRAOne 단위)에 무너진다. 배포한 컨트랙트마다 자기 몫의 재시작 질량을 가져오기 때문이다. 균등 재시작 EndorseRank는 두 공격 모두에 무너진다. 명제~\ref{prop:bipartite}가 예측하듯 지갑 \ToyCostERUniformATwo 개면 충분해서 호출자 수 세기에 필요한 \ToyCostCallersATwo 개와 큰 차이가 없고, 농장 컨트랙트도 재시작 질량을 가진 노드이므로 통한다. \PCER는 명제~\ref{prop:region}에서 보듯 A1에는 크기와 관계없이 움직이지 않고, A2에는 \ToyCostPCERATwo 단위가 들어 균등 재시작일 때의 약 \ToyRatioATwo 배이다. 이 배수는 \ref{sec:method-weights}절의 예측대로 정직한 지갑이 낸 평균 수수료(최소 수수료의 \ToyMeanFee 배)에 가깝다. A4는 실행 시점 trait 바인딩을 라우터 간선으로 더했을 때만 토큰을 움직인다(\ToyCostPCERBindAFour 단위). 이 그래프에서는 이 공격이 A2보다 비싸다. 정직한 호출이 이미 라우터의 토큰을 수백 번 바인딩하고 있기 때문이다. 점수가 호출보다 위임에서 주로 오는 라우터라면 더 싸게 장악할 수 있다.

\subsection{클론 묶음, 패키지, 검색}
\label{sec:method-search}

\textbf{클론 묶음.} 정규화한 소스의 해시가 같은 두 컨트랙트는 같은 묶음에 속한다. 정규화는 주석과 공백을 지우고, 문자열·숫자 리터럴, principal 리터럴, 정의한 토큰·맵·변수의 이름을 자리표시자로 바꾼 뒤 토큰 열을 해시한다. 코드 본문이 똑같은 컨트랙트들은 \code{contract-hash?} 값도 같다. MinHash \cite{broder1997resemblance}로 찾은 Jaccard 유사도 0.9 이상의 근사 중복은 관련 컨트랙트로 보여 주되 합치지는 않는다. 묶음은 점수가 가장 높은 구성원이 대표하고(동점이면 먼저 배포된 것), 묶음의 순위는 그 구성원의 점수로 정한다. 합 대신 최댓값을 쓰므로, 어떤 컨트랙트의 복제본이 천 개 있어도 그 묶음에는 아무것도 보태지 않는다.

\textbf{패키지.} 패키지는 게시자가 선언한 컨트랙트의 계보이다(\ref{sec:registry}절). 패키지의 순위는 최신 버전의 점수로 정하고, 이전 버전은 그 아래에 나열한다.

\textbf{검색.} 식~\eqref{eq:pkgsite}를 따라, 질의 $q$와 맞는 컨트랙트 $c$에 다음 점수를 준다.
\begin{equation}
\label{eq:search}
\mathrm{score}(q,c)=\mathrm{BM25}(q,c)\cdot\bigl(\ln(e+|C|\,r_c)\bigr)^{\alpha},\qquad \alpha=1
\end{equation}
여기서 BM25 \cite{robertson2009bm25}는 컨트랙트 이름, 게시자 스코프, ABI의 함수 이름과 trait 이름, 소스의 주석, 패키지 README를 대상으로 계산하고, $|C|\,r_c$는 점수가 평균인 컨트랙트에서 1이다. 결과는 trait 준수 여부(SIP-009, SIP-010, SIP-013), Clarity 버전, 네트워크로 거를 수 있다.
